%% P26 --- Estymacja i metoda najwiekszej wiarogodnosci
%%
%% Blizniak: programs/en/P26-maximum-likelihood.tex. Ramka w ramke, makro w
%% makro, liczba w liczbe; tools/parity.py porownuje oba pliki. Kazda liczba
%% pochodzi z code/p26_maximum_likelihood.py.
%%
%% UWAGI DLA TLUMACZA, z regul zapisanych w CLAUDE.md:
%%   * cyfra zostaje cyfra ORAZ slowo zostaje slowem -- P25 zlapal ten blad
%%     w druga strone: polski napisal $2$ i $4$ tam, gdzie angielski pisze
%%     "two" i "four", i C4, C8 oraz C12 odpalily naraz;
%%   * odwolanie prowadzi, a nie idzie za matematyka;
%%   * gdzie angielski pisze skrot, polski tez musi;
%%   * terminy ML zostaja angielskie i odmieniaja sie po polsku.

\program[Estymacja i wiarogodność]{Estymacja i metoda największej
wiarogodności}
\label{prog:P26}
\index{estymator}
\index{największa wiarogodność}

\begin{outcomes}
\outcome{1--8}{nazwać dwa sposoby, na jakie estymator może być zły, i
  rozłożyć na nie błąd średniokwadratowy}
\outcome{10--17}{powiedzieć, co poprawia $n-1$, a czego nie poprawia}
\outcome{18--23}{zapisać wiarogodność i znaleźć parametr, który ją
  maksymalizuje}
\outcome{25--33}{wyprowadzić stratę treningową modelu językowego z
  prawdopodobieństwa jego danych}
\outcome{34--39}{odczytać weight decay jako stwierdzenie o tym, jak duże
  według ciebie są wagi}
\outcome{41--46}{zapisać score i użyć go do różniczkowania przez losowanie}
\end{outcomes}

\begin{quiz}
\item Estymator jest nieobciążony. Podaj coś, co i tak czyni go bezużytecznym.
  \teachesat{1--2}
  \answerto{Wariancja. Nieobciążoność mówi, że zbiór ocen jest wyśrodkowany
  na prawdzie; nie mówi nic o tym, jak daleko od niej leży którakolwiek z
  nich, a to są dwie osobne osie.}
\item Wariancja z próby dzieli przez $n-1$, a nie przez $n$. Co dokładnie to
  naprawia? \teachesat{11--12}
  \answerto{Jedno obciążenie, dokładnie: wersja z mianownikiem $n$ daje
  średnio $\frac{n-1}{n}$ wariancji populacji, a dzielenie przez $n-1$
  sprawia, że średnia wychodzi dokładnie taka, jaka być powinna.}
\item A czy pierwiastek z tej poprawionej wariancji jest nieobciążoną oceną
  odchylenia standardowego? \teachesat{13--14}
  \answerto{Nie. Pierwiastek jest wklęsły, więc z nierówności Jensena jego
  średnia leży poniżej pierwiastka ze średniej \dash{} poprawione odchylenie
  standardowe z próby zaniża, a \code{ddof=1} tego nie naprawia.}
\item Model przypisuje prawdopodobieństwa tokenom, które zobaczył. Co robi z
  tymi prawdopodobieństwami trenowanie go na minimalizację entropii
  krzyżowej? \teachesat{26--27}
  \answerto{Maksymalizuje ich iloczyn. Entropia krzyżowa wobec celu one-hot
  jest ujemnym logarytmem wiarogodności zaobserwowanych tokenów podzielonym
  przez ich liczbę, więc jej minimalizacja to metoda największej
  wiarogodności i nic więcej.}
\item Dodajesz weight decay o sile $\lambda$. Powiedz, co to mówi o wagach, a
  nie co im robi. \teachesat{37--38}
  \answerto{Że według ciebie są gaussowskie wokół zera o pewnej szerokości.
  Kara jest logarytmem rozkładu a priori, a szerokość ustala $\lambda$.}
\end{quiz}

% ============================================================
\section{Dwa sposoby bycia w błędzie}
\label{sec:P26-two-ways}
\index{estymator!obciążenie i wariancja}
\index{pomiar!błąd średniokwadratowy}

\begin{fr}
Program~\ref{prog:P21} rozebrał gradient z paczki i ustalił o nim jedną
rzecz: uśredniony po każdej paczce, jaka mogła zostać wylosowana, jest tym
gradientem, o który ci chodziło. To znaczy \emph{nieobciążony} i program
dowiódł tego dokładnie, a nie przez pokaz.

Potem jego własna ramka o pułapce powiedziała, że nieobciążoność nie
wystarcza. Ten program jest o tym, czego nie wystarcza.

Zanim pójdziesz dalej: ocena jest nieobciążona. Podaj drugi sposób, na jaki
i tak może być zła.
\dotline
\nextframe
\end{fr}

\begin{fr}
\begin{ansblock}
Może być rozrzucona na wszystkie strony. Nieobciążoność jest stwierdzeniem o
\emph{zbiorze} ocen i nie mówi nic o żadnej pojedynczej.
\end{ansblock}

Są więc dwie osie i są od siebie niezależne.

\textbf{Obciążenie} to odległość średniej oceny od prawdy.
\textbf{Wariancja} to odległość ocen od ich własnej średniej. Estymator może
być zły na jeden z tych sposobów, na oba albo \dash{} pewnym kosztem \dash{}
na jeden, żeby kupić drugi.

Wszystko w tej sekcji jest dokładne. Weź populację
$\val{p26.pop.n}$ liczb $1$, $2$, $3$, $6$, losuj
$\val{p26.draw.n}$ niezależnie, a możliwych prób jest $\val{p26.samples}$;
skrypt wylicza je wszystkie po kolei nad ułamkami i nie ma w tym nigdzie
żadnej tolerancji.

Średnia populacji to $\val{p26.pop.mean}$, a jej wariancja to
$\val{p26.pop.var}$. Jakie są obciążenie i wariancja średniej z próby?
\dotline
\nextframe
\end{fr}

\begin{fr}
\ans{Obciążenie $0$, wariancja $\val{p26.var.unbiased}$}

Nieobciążona, co jest wynikiem Programu~\ref{prog:P21} na tej populacji, i
wariancja $\sigma^{2}/n$, co jest tempem Programu~\ref{prog:P25}. Nic tu
nowego; to punkt odniesienia, względem którego mierzy się reszta sekcji.

Teraz zmień estymator. Zamiast średniej z próby $\bar{x}$ weź
$c\,\bar{x}$ dla ustalonego $c$ między $0$ a $1$ \dash{} ta sama ocena,
przyciągnięta ku zeru.

Powiedz, co to robi z każdą z dwóch osi.
\dotline
\nextframe
\end{fr}

\begin{fr}
\begin{ansblock}
Wprowadza obciążenie i mnoży wariancję przez $c^{2}$.
\end{ansblock}

Obciążenie wynosi $(c-1)\mu$, bo średnia ocena to teraz $c\mu$, a nie $\mu$.
Wariancja niesie kwadrat, bo Program~\ref{prog:P24} podaje
$\Var(cX) = c^{2}\Var(X)$.

Przyciąganie wymienia więc jedną oś na drugą i oczywiste pytanie brzmi, czy
ta wymiana kiedykolwiek się opłaca. Żeby je w ogóle zadać, potrzebujesz
jednej liczby, do której wchodzą obie osie.

Co byś zmierzył?
\dotline
\nextframe
\end{fr}

\begin{fr}
\begin{ansblock}
Średnią kwadratową odległość od prawdy \dash{} \textbf{błąd
średniokwadratowy}.
\end{ansblock}

I on się rozkłada, dokładnie:
\[ \mathrm{MSE} \;=\; \Ex\bigl[(\hat\theta - \theta)^{2}\bigr]
   \;=\; \text{obciążenie}^{2} + \Var(\hat\theta) \]

Skrypt sprawdza to jako tożsamość przy dwudziestu jeden przyciąganiach, a nie
przy jednej wygodnej wartości, nad ułamkami \dash{} więc jest to fakt o
estymatorze, a nie arytmetyczny przypadek.

Kwadrat odległości robi w tym zdaniu robotę i warto powiedzieć dlaczego.
Średni błąd bezwzględny nie rozkłada się wcale: nie ma tożsamości, która
rozdzieliłaby go na część o środku i część o rozrzucie. To podniesienie do
kwadratu sprawia, że obie osie się dodają, zamiast oddziaływać, i z tego
samego powodu Program~\ref{prog:P24} trzyma wariancję, a nie średnią
odległość; dlatego właśnie to ta strata, a nie inna, ma rozkład, o którym da
się mówić.

Przeczytaj to, zanim przewrócisz stronę: czy $c = 1$ \dash{} wybór
nieobciążony \dash{} minimalizuje ten błąd?
\dotline
\nextframe
\end{fr}

\begin{fr}
\begin{ansblock}
Nie. Na tej populacji najlepsze $c$ to $\val{p26.shrink.c}$.
\end{ansblock}

\begin{center}
\begin{tabular}{lrr}
\toprule
& $c = 1$ & $c = \val{p26.shrink.c}$ \\
\midrule
obciążenie & $0$ & $\val{p26.shrink.bias}$ \\
wariancja & $\val{p26.var.unbiased}$ & $\val{p26.shrink.var}$ \\
błąd średniokwadratowy & $\val{p26.mse.unbiased}$ & $\val{p26.mse.shrunk}$ \\
\bottomrule
\end{tabular}
\end{center}

\textbf{Celowo obciążony estymator jest bliżej prawdy}, o
$\val{p26.shrink.gain.pct}$ procent błędu średniokwadratowego, a obie kolumny
sumują się do własnego dolnego wiersza co do wydrukowanej cyfry.
\end{fr}

\begin{fr}
\begin{trapbox}
\textbf{\enquote{Nieobciążoność to jest to, o co chodzi.}}

Nieobciążoność jest jedną z dwóch osi i tą, którą łatwo sprawdzić, dlatego
jest cytowana. Nie jest tą, która mówi, jak daleko twoja odpowiedź leży od
prawdy. Nieobciążony estymator o dużej wariancji jest gorszy, według każdej
miary, na jakiej praktykowi zależy, od lekko obciążonego, który stoi
spokojnie.

Podejrzliwość należy się słowu \emph{najlepszy}. Najlepszy w czym? Nic nie
jest najlepsze, dopóki nie nazwano straty, a \enquote{nieobciążony} nie jest
stratą.
\end{trapbox}

Jedno uczciwe zastrzeżenie, bo to ono nie pozwala z tego zrobić przepisu:
optymalne $c$ powyżej zależy od $\mu$, czyli od tej wielkości, którą się
szacuje. Nie da się go policzyć.

Co jest więc warte to spostrzeżenie, skoro najlepsze przyciąganie jest
niedostępne?
\dotline
\nextframe
\end{fr}

\begin{fr}
\begin{ansblock}
To, że obciążenie wprowadzone celowo bywa dobrą wymianą \dash{} czym jest
każda regularyzacja w tej książce.
\end{ansblock}

Jedną już spotkałeś. Program~\ref{prog:P20} podaje weight decay, które
przyciąga każdy parametr ku zeru, obciążając każdy z nich, i używa się go
dlatego, że wariancja, którą usuwa, jest warta więcej niż obciążenie, które
dodaje.

To ta sama wymiana co w tabeli powyżej, na modelu zamiast na średniej, a
\S\ref{sec:P26-prior} mówi dokładnie, czym jest w niej $\lambda$.
\end{fr}

\begin{fr}
\mermaidfig{p26-two-axes}{Dwie osie i jedna liczba, która je
dodaje.}{obciazenie, wariancja, blad}

Dwa sposoby bycia w błędzie i strata, która je waży. Reszta programu jest o
tym, skąd w ogóle bierze się estymator.
\end{fr}

% ============================================================
\section{$n-1$ i to, czego ono nie naprawia}
\label{sec:P26-n-minus-one}
\index{estymator!poprawka n-1}
\index{wariancja!z próby}

\begin{fr}
Najsłynniejsze obciążenie w statystyce i takie, którego prawie nikt nie umie
wyjaśnić.

Wariancja z próby dzieli przez $n-1$, a nie przez $n$. Każda biblioteka ma na
to flagę i każdy kurs to podaje. Zanim pójdziesz dalej, powiedz, co idzie
źle, gdy podzielisz przez $n$.
\dotline
\nextframe
\end{fr}

\begin{fr}
\begin{ansblock}
Wychodzi za mała, i to nie czasem \dash{} średnio, za każdym razem.
\end{ansblock}

Powód to jedno zdanie i warto je mieć, a nie pamiętać. Wariancja z próby
mierzy rozrzut wokół średniej \emph{z próby}, a średnia z próby jest
dopasowana do tej właśnie próby, więc leży bliżej danych niż prawdziwa
średnia. Odległości mierzone od niej są zatem za krótkie.

Zauważ, od czego ten argument nie zależy. Nie mówi nic o tym, które liczby
zostały wylosowane, a tylko o tym, jak zbudowano estymator, więc niedobór
jest własnością przepisu, a nie twoich danych. Dlatego poprawka może być
stałym czynnikiem stosowanym raz, a nie czymś szacowanym obok wszystkiego
innego.

To jest obciążenie, więc ma swój rozmiar. Zgadnij czynnik, zanim przewrócisz
stronę.
\dotline
\nextframe
\end{fr}

\begin{fr}
\ans{Dokładnie $\frac{n-1}{n}$}

Nie w przybliżeniu i nie tylko dla dużych $n$. Wyliczone po każdej próbie
przy $n = 2$, $3$, $4$ i $5$ na populacji z
\S\ref{sec:P26-two-ways}, nad ułamkami:

\begin{center}
\begin{tabular}{lr}
\toprule
wariancja populacji & $\val{p26.pop.var}$ \\
średnia wariancji z próby z mianownikiem $n$ &
  $\val{p26.var.n.expected}$ \\
średnia wariancji z próby z mianownikiem $(n-1)$ &
  $\val{p26.var.n1.expected}$ \\
\bottomrule
\end{tabular}
\end{center}

Przy $n = \val{p26.draw.n}$ naiwna wersja jest za mała o
$\val{p26.var.understate.pct}$ procent, bo $\frac{n-1}{n}$ wynosi tam połowę.
Dzielenie przez $n-1$ sprawia, że średnia wychodzi dokładnie taka, jaka
powinna, i na tym polega całość tego, co ta poprawka poprawia.

To jest więc załatwione. Teraz pytanie, do którego zachęca nazwa flagi: czy
pierwiastek z poprawionej wariancji jest nieobciążoną oceną odchylenia
standardowego?
\dotline
\nextframe
\end{fr}

\begin{fr}
\begin{ansblock}
Nie, i to bardziej, niż byś się spodziewał.
\end{ansblock}

Ta sama populacja, to samo wyliczenie, poprawione odchylenie standardowe z
próby uśrednione po każdej próbie:

\begin{center}
\begin{tabular}{rrr}
\toprule
$n$ & średnie poprawione $s$ & za małe o \\
\midrule
$2$ & $\val{p26.sd.mean.2}$ & $\val{p26.sd.short.2}\%$ \\
$3$ & $\val{p26.sd.mean.3}$ & $\val{p26.sd.short.3}\%$ \\
$5$ & $\val{p26.sd.mean.5}$ & $\val{p26.sd.short.5}\%$ \\
$10$ & $\val{p26.sd.mean.10}$ & $\val{p26.sd.short.10}\%$ \\
\bottomrule
\end{tabular}
\end{center}

Odchylenie standardowe populacji wynosi $\val{p26.pop.sd}$. Przy $n = 2$
poprawiona ocena jest za mała o niemal jedną czwartą.

Program~\ref{prog:P19} podaje powód w jednym wierszu. Powiedz, w którym.
\dotline
\nextframe
\end{fr}

\begin{fr}
\begin{ansblock}
Nierówność Jensena, dla funkcji \textbf{wklęsłej}.
\end{ansblock}

Pierwiastek jest wklęsły, a Program~\ref{prog:P19} podaje przypadek wklęsły w
tylu właśnie słowach: funkcja średniej jest co najmniej średnią funkcji, więc
\[ \Ex\bigl[\sqrt{s^{2}}\bigr] \;\le\; \sqrt{\Ex[s^{2}]} \;=\; \sigma \]

Nieobciążona ocena wielkości nie przeżywa nałożenia na nią funkcji
nieliniowej. \textbf{Poprawienie wariancji nie może poprawić odchylenia
standardowego, bo poprawka i pierwiastek nie są przemienne.}

A przypadek równości z Programu~\ref{prog:P19} mówi, kiedy luka się domyka:
tylko wtedy, gdy wariancja z próby jest ta sama dla każdej próby, a nie jest.
Rozbieżność powyżej maleje z $n$ z tego samego powodu \dash{} oceny mniej się
rozrzucają.
\end{fr}

\begin{fr}
\begin{trapbox}
\textbf{\enquote{\code{ddof=1} daje mi nieobciążone odchylenie standardowe.}}

Daje nieobciążoną \emph{wariancję}. Żadna biblioteka nie daje nieobciążonego
odchylenia standardowego, bo poprawka zależy od rozkładu i nie ma na nią
flagi.

Wart wiedzy, a nie pedanterii, czyni to rozmiar. Przy pięciu przebiegach
eksperymentu \dash{} a to właśnie od tylu ludzie raportują odchylenie
standardowe \dash{} liczba jest za mała o jakieś
$\val{p26.sd.short.5}$ procent, i to za mała w tę stronę, która sprawia, że
wyniki wyglądają na powtarzalniejsze, niż są.
\end{trapbox}
\end{fr}

\begin{fr}
\transcript{p26-still-short}

Dwa porównania. Pierwsze mówi, że poprawiona wariancja jest dobra, a naiwna
nie; drugie drukuje poprawione odchylenie standardowe obok tego z populacji.
\end{fr}

\begin{fr}
Co więc należy z tym zrobić?

Nic dramatycznego, a uczciwa odpowiedź ma trzy części. Podawaj wariancję,
skoro czytelnik może sam wyciągnąć pierwiastek, bo to wariancja jest tą
wielkością, która się dodaje, i to dlatego ta książka ją trzyma. Znaj rozmiar
obciążenia przy tym $n$, które masz. I traktuj odchylenie standardowe z
kilku przebiegów jako przybliżoną liczbę, którą ono jest.

Nic z tego nie wymaga czynnika poprawkowego, a sięgnięcie po niego jest
sposobem, w jaki przemyca się założenie o rozkładzie.
\end{fr}

% ============================================================
\section{Parametr, przy którym dane najmniej zaskakują}
\label{sec:P26-mle}
\index{największa wiarogodność!definicja}
\index{wiarogodność!jako funkcja parametru}

\begin{fr}
Obie dotychczasowe sekcje brały estymator jako dany i pytały, jak bardzo jest
zły. Teraz pytanie leżące pod spodem: skąd w ogóle bierze się estymator?

Masz dane i rodzinę rozkładów indeksowaną parametrem. Powiedz, co czyniłoby
jedną wartość tego parametru lepszą od innej.
\dotline
\nextframe
\end{fr}

\begin{fr}
\begin{ansblock}
To, że czyni bardziej prawdopodobnymi te dane, które naprawdę zobaczyłeś.
\end{ansblock}

Na tym polega cały pomysł i jest to reguła zamiany danych w liczbę, a nie
fakt o świecie. \textbf{Wiarogodność} to
$\Prob(\text{dane} \mid \theta)$, czytane jako funkcja $\theta$ przy
ustalonych danych, a \textbf{metoda największej wiarogodności} wybiera to
$\theta$, które ją maksymalizuje.

Program~\ref{prog:P23} już ostrzegał przed tym słowem i to ostrzeżenie ma tu
znaczenie: wiarogodność jest prawdopodobieństwem \emph{danych}, a nie
parametru. Nie całkuje się do jedynki po $\theta$ i nie jest rozkładem na
$\theta$. Traktowaniem jej jak rozkładu zajmuje się porządnie
Program~\ref{prog:P28}, z rozkładem a priori, a ta sekcja nie.

Dane to $n$ niezależnych obserwacji, więc wiarogodność jest iloczynem $n$
czynników. Podaj dwa powody, by wziąć jej logarytm.
\dotline
\nextframe
\end{fr}

\begin{fr}
\begin{ansblock}
Zamienia iloczyn w sumę i powstrzymuje iloczyn przed niedomiarem.
\end{ansblock}

Pierwszy należy do Programu~\ref{prog:F03}, czyli do miejsca, w którym ta
książka zamienia iloczyny w sumy, i to on czyni pochodną wykonalną: suma
różniczkuje się składnik po składniku, a iloczyn $n$ rzeczy nie. Drugi należy
do Programu~\ref{prog:P02} \dash{} własne prawdopodobieństwo sekwencji z
Programu~\ref{prog:F03} wynosiło $\val{f03.seq.prod}$, czyli liczbę, której
żaden format nie utrzyma.

A logarytm nie może przesunąć maksimum, bo Program~\ref{prog:F05} dowiódł, że
funkcja ściśle rosnąca zachowuje porządek listy.

Zatem: zapisz logarytm wiarogodności, zróżniczkuj, przyrównaj do zera, co
jest przepisem Programu~\ref{prog:P15} na maksimum. Spróbuj. Dwadzieścia prób
daje siedem sukcesów; jaka częstość maksymalizuje wiarogodność?
\dotline
\nextframe
\end{fr}

\begin{fr}
\ans{$\val{p26.bern.mle}$}

Oczywista odpowiedź i warto zobaczyć, że jest wyprowadzona, a nie założona.
Przy $k$ sukcesach w $n$ próbach logarytm wiarogodności to
$k\ln p + (n-k)\ln(1-p)$ z dokładnością do stałej, a
\[ \frac{\mathrm{d}}{\mathrm{d}p}
   \Bigl[k\ln p + (n-k)\ln(1-p)\Bigr]
   \;=\; \frac{k}{p} - \frac{n-k}{1-p} \]
zeruje się dokładnie przy $p = k/n$. Skrypt sprawdza to nad ułamkami, a potem
sprawdza, że każda inna częstość na siatce tysiąca punktów jest ściśle
gorsza, więc jest to maksimum, a nie punkt stacjonarny.

\textbf{Każda częstość, jaką ktokolwiek policzył, dzieląc sukcesy przez
próby, była oceną największej wiarogodności}, i książka mówi to po raz
pierwszy.
\end{fr}

\begin{fr}
Teraz ten sam przepis na rozkładzie Gaussa, bo jego odpowiedź zderza się z
\S\ref{sec:P26-n-minus-one} celowo.

Maksymalizacja logarytmu wiarogodności $\val{p26.gauss.n}$ obserwacji po
średniej daje średnią z próby, tutaj $\val{p26.gauss.mu}$. Po wariancji daje
średnią kwadratową odległość od tej średniej \dash{} z $n$ w mianowniku, co
wychodzi $\val{p26.gauss.mle.var}$ wobec nieobciążonej
$\val{p26.gauss.unbiased.var}$.

Powiedz, czym to czyni metodę największej wiarogodności.
\dotline
\nextframe
\end{fr}

\begin{fr}
\begin{ansblock}
Obciążoną.
\end{ansblock}

\begin{trapbox}
\textbf{\enquote{Metoda największej wiarogodności daje właściwą odpowiedź.}}

Daje tę odpowiedź, która czyni twoje dane najbardziej prawdopodobnymi, co
jest czym innym i nie znaczy nieobciążoną. Wariancja gaussowska jest tu
przykładem wzorcowym: ocena, którą metoda największej wiarogodności ci wręcza,
jest za mała dokładnie o $\frac{n-1}{n}$ z poprzedniej sekcji, a poprawka,
którą wszyscy stosują, jest poprawką \emph{od} odpowiedzi największej
wiarogodności.

Dwie metody, dwie odpowiedzi, a ta niezgoda nie jest błędem żadnej z nich.
Optymalizują różne rzeczy.
\end{trapbox}
\end{fr}

\begin{fr}
\begin{rigourbox}
\textbf{Nie dowodzimy tutaj: dlaczego w ogóle warto używać metody największej
wiarogodności.}

Argument za nią jest asymptotyczny. Gdy $n$ rośnie, ocena zbiega do
prawdziwego parametru, a wśród szerokiej klasy estymatorów jej wariancja
dąży do najmniejszej osiągalnej. Oba stwierdzenia są twierdzeniami z
założeniami, oba potrzebują maszynerii spoza tej książki, a drugie ma nazwę
\dash{} nierówność Craméra--Rao \dash{} którą czytelnik spotka i może
sprawdzić.

Warto natomiast zabrać kształt tej obietnicy. Jest \emph{asymptotyczna}:
metoda największej wiarogodności jest dobrym pomysłem w granicy i jest
obciążona przy każdym skończonym $n$, co dokładnie pokazuje wariancja
gaussowska. Nic w niej nie mówi, że odpowiedź przy twoim rozmiarze próby jest
nieobciążona, a poprzednia sekcja mówi, że nie jest.
\end{rigourbox}
\end{fr}

% ============================================================
\section{Trenowanie modelu językowego to metoda największej wiarogodności}
\label{sec:P26-training}
\index{entropia krzyżowa!jako ujemny logarytm wiarogodności}
\index{największa wiarogodność!a modele językowe}

\begin{fr}
Po to była poprzednia sekcja i to jest ta ramka, którą obiecał
Program~\ref{prog:P18}.

Model czyta sekwencję i na każdej pozycji produkuje rozkład na słowniku.
Token, który naprawdę wystąpił jako następny, ma pod tym rozkładem pewne
prawdopodobieństwo. Nazwij parametry modelu $\theta$, a zaobserwowane tokeny
$t_1, \dots, t_n$.

Zapisz prawdopodobieństwo, jakie model przypisuje całej zaobserwowanej
sekwencji.
\dotline
\nextframe
\end{fr}

\begin{fr}
\ans{$\prod_{i=1}^{n} p_\theta(t_i \mid t_{<i})$}

Co jest wiarogodnością: prawdopodobieństwem danych czytanym jako funkcja
parametrów. Metoda największej wiarogodności każe więc maksymalizować to po
$\theta$.

Zrób teraz te dwie rzeczy, które kazała robić \S\ref{sec:P26-mle}. Weź
logarytm, który zamienia to w sumę. Potem, skoro optymalizatory minimalizują,
zmień znak. Potem podziel przez $n$, żeby sekwencje różnej długości były
porównywalne.

Zmiana znaku jest jedyną z tych trzech rzeczy, która jest czystą konwencją.
Maksymalizowanie wielkości i minimalizowanie jej wartości przeciwnej to ten
sam problem, a każdy optymalizator z Programu~\ref{prog:P20} był pisany tak,
by schodzić w dół, więc znak jest tam dla maszynerii, a nie dla matematyki. W
zamian kupuje to, że liczba czyta się jak koszt, i dlatego wynik jednego
przebiegu jest porównywalny z wynikiem innego.

Zapisz, co ci wyszło.
\dotline
\nextframe
\end{fr}

\begin{fr}
\ans{$-\frac{1}{n}\sum_{i=1}^{n} \ln p_\theta(t_i \mid t_{<i})$}

\begin{note}
\textbf{Definicja, którą Program~\ref{prog:P18} podał i zostawił bez
uzasadnienia.} Tamten program potrzebował pochodnej entropii krzyżowej
znacznie wcześniej, niż ta książka mogła powiedzieć, po co ta strata jest,
więc zdefiniował $L = -\sum_j y_j \ln p_j$, powiedział, że uzasadnienie jest
tutaj, i zajął się jakobianem.

Oto ono. Przy celu one-hot wszystkie składniki poza jednym są zerowe i
$L = -\ln p_c$ dla zaobserwowanego tokenu, więc średnia entropia krzyżowa po
sekwencji \textbf{jest} wyrażeniem powyżej, znak po znaku.

\textbf{Strata krzyżowoentropijna jest ujemnym logarytmem wiarogodności
zaobserwowanych tokenów, na token, i nie ma w niej nic więcej.} Trenowanie
modelu językowego to metoda największej wiarogodności. Nigdy nie było niczym
innym.
\end{note}
\end{fr}

\begin{fr}
Warto prześledzić ten łańcuch wstecz, bo okazuje się, że trzy programy cytują
jedno obliczenie.

Program~\ref{prog:F02} podaje stratę $\val{p26.f02.loss}$ nata na token.
Program~\ref{prog:F03} bierze ją, zamienia w prawdopodobieństwo na token i
podnosi je do potęgi $\val{p26.f03.tokens}$ tokenów, dostając
prawdopodobieństwo sekwencji, którego logarytm dziesiętny zapisuje.

Odwróć to: weź ten zapisany logarytm, podziel przez liczbę tokenów i zmień
znak. Skrypt robi dokładnie to i dostaje $\val{p26.back.loss}$ nata, zgodnie z
Programem~\ref{prog:F02} z dokładnością do tego, co może wyjaśnić
zaokrąglenie zapisanego wykładnika do dwóch miejsc \dash{} tolerancja
wyprowadzona, a nie wybrana.

\textbf{Strata i prawdopodobieństwo sekwencji to ta sama liczba zapisana na
dwa sposoby}, a budowanie przerywa teraz, jeśli któryś z programów się
ruszy.
\end{fr}

\begin{fr}
\mermaidfig{p26-one-quantity}{Jedna wielkość, do której dochodzi się z dwóch
stron.}{wiarogodnosc i strata}

Z tego przeramowania wypadają trzy konsekwencje i każda dotyczy czegoś, co
ludzie robią, nie mając pod ręką powodu.
\end{fr}

\begin{fr}
Pierwsza. \textbf{Label smoothing} zastępuje cel one-hot wartością
$1-\varepsilon$ na zaobserwowanym tokenie i $\varepsilon/K$ rozłożonym po
wszystkich $K$ z nich, i zwykle wprowadza się je jako sztuczkę
regularyzacyjną.

Przeczytaj je zamiast tego jako metodę największej wiarogodności. Strata
$-\sum_j t_j \ln p_j$ jest minimalizowana na sympleksie wtedy, gdy $\vect{p}$
równa się $\vect{t}$, jakiekolwiek $\vect{t}$ by nie było.

Dla jakiego więc celu label smoothing jest metodą największej wiarogodności?
\dotline
\nextframe
\end{fr}

\begin{fr}
\begin{ansblock}
Dla rozkładu wygładzonego \dash{} celu, którego nigdy nie zaobserwowano.
\end{ansblock}

A to jest prawdziwe stwierdzenie o tym, co ono robi. Przy
$\varepsilon = \val{p26.ls.eps}$ na $\val{p26.ls.classes}$ klasach optimum to
$\val{p26.ls.true}$ na zaobserwowanym tokenie i $\val{p26.ls.other}$ na
każdym z pozostałych \dash{} znalezione przez przeszukanie, a nie zacytowane
\dash{} więc optymalna różnica między zwycięskim logitem a resztą wynosi
$\val{p26.ls.gap}$.

\textbf{Na tym to polega i jest to liczba.} Wobec celu one-hot optimum chce
$p_c = 1$, co wymaga nieskończonej różnicy logitów, więc strata popycha bez
końca i nigdy nie dociera. Wygładzanie nazywa zamiast tego optimum
osiągalne. To nie jest szum dodany na szczęście; to inny, osiągalny cel.
\end{fr}

\begin{fr}
Druga. \textbf{Ważenie klas} mnoży stratę każdego przykładu przez wagę i
czyta się jako \enquote{bardziej mi zależy na klasie rzadkiej}.

Jako metoda największej wiarogodności jest ostrzejsze: zważenie przykładu
przez $w$ to logarytm wiarogodności próby, w której ten przykład wystąpił $w$
razy. Ważenie nie czyni więc modelu lepszym na klasie rzadkiej \dash{} ono
szacuje parametry \emph{innej populacji}, takiej, w której klasa rzadka jest
częstsza niż w twojej.

Czy o to ci chodzi, jest pytaniem o wdrożenie, a to przeramowanie pozwala je
zadać.

Trzecia jest mniejsza i już jest na stronie. Powiedz, dlaczego strata jest
dzielona przez liczbę tokenów, a nie sumowana.
\dotline
\nextframe
\end{fr}

\begin{fr}
\begin{ansblock}
Żeby długa i krótka sekwencja były porównywalne.
\end{ansblock}

Logarytm wiarogodności rośnie z liczbą tokenów, więc suma mierzy korpus tak
samo jak model. Na token jest średnią, a Program~\ref{prog:P19} mówi, co
uśrednianie straty na token przeżywa, a czego nie \dash{} wolno ci uśredniać
straty i nie wolno ci uśredniać perplexity.

\begin{trapbox}
\textbf{\enquote{Entropia krzyżowa mierzy odległość między dwoma
rozkładami.}}

Nie jest odległością: nie znika, gdy oba się zgadzają, i nie jest
symetryczna. Znika \emph{nadwyżka} entropii krzyżowej nad własną entropią
celu, a ta nadwyżka ma nazwę i należy do Programu~\ref{prog:P30}.

Dla tego programu odczyt jest węższy i kompletny: wobec celu one-hot entropia
krzyżowa jest ujemnym logarytmem wiarogodności, a jej minimum leży tam, gdzie
wiarogodność jest największa.
\end{trapbox}
\end{fr}

% ============================================================
\section{Kara jest logarytmem rozkładu a priori}
\label{sec:P26-prior}
\index{regularyzacja!jako rozkład a priori}
\index{weight decay!czym jest jego lambda}

\begin{fr}
Metoda największej wiarogodności pyta tylko o to, co czyni dane
prawdopodobnymi. Przypuśćmy, że masz też zdanie o parametrze jeszcze przed
zobaczeniem danych \dash{} że wagi sieci są raczej małe, powiedzmy.

Program~\ref{prog:P23} daje maszynerię: pomnóż wiarogodność przez rozkład a
priori, a dostaniesz coś proporcjonalnego do rozkładu a posteriori.
Maksymalizowanie tego zamiast tamtego to \textbf{maksimum a posteriori},
czyli MAP.

Weź logarytm iloczynu dwóch rzeczy. Co dzieje się z funkcją celu?
\dotline
\nextframe
\end{fr}

\begin{fr}
\begin{ansblock}
Logarytm rozkładu a priori zostaje \emph{dodany} do logarytmu wiarogodności.
\end{ansblock}

\[ \ln\bigl[\Prob(\text{dane} \mid \theta)\,\Prob(\theta)\bigr]
   \;=\; \ln \Prob(\text{dane} \mid \theta) + \ln \Prob(\theta) \]

Zmień znak, jak zrobiła to \S\ref{sec:P26-training}, a drugi składnik staje
się czymś dodanym do straty i minimalizowanym razem z nią.

Taki jest kształt każdej regularyzowanej funkcji celu, jaką ktokolwiek
zapisał. Nałóż teraz na każdą wagę gaussowski rozkład a priori o szerokości
$\tau$, wyśrodkowany na zerze, i powiedz, jaki jest dodany składnik.
\dotline
\nextframe
\end{fr}

\begin{fr}
\ans{$\dfrac{\lVert \vect{w}\rVert^{2}}{2\tau^{2}}$, plus stała}

Bo logarytm $e^{-w^{2}/2\tau^{2}}$ to $-w^{2}/2\tau^{2}$, a czynnik
normalizujący nie zależy od $\vect{w}$, więc nie może przesunąć minimum.

\textbf{Kara kwadratowa jest logarytmem gaussowskiego rozkładu a priori.} Na
tym polega cała ta odpowiedniość i jest to jeden wiersz.

Kara nie jest więc urządzeniem doczepionym do funkcji celu na szczęście. Jest
tym, co dostajesz, zapisując przekonanie i biorąc logarytm, a siła kary jest
stwierdzeniem o \emph{szerokości} tego przekonania.
\end{fr}

\begin{fr}
Co pozwala odczytać hiperparametr, zamiast go stroić.

Program~\ref{prog:P20} rozstrzygnął weight decay, dodając $\lambda\vect{w}$
do gradientu, i zmierzył, co to robi. Dopasuj oba gradienty \dash{}
$\lambda\vect{w}$ wobec $\vect{w}/\tau^{2}$ \dash{} a wyjdzie
\[ \lambda \;=\; \frac{1}{\tau^{2}}
   \qquad\text{czyli}\qquad \tau \;=\; \frac{1}{\sqrt{\lambda}} \]

Skrypt dopasowuje je na tych gradientach, które Program~\ref{prog:P20}
naprawdę dodaje, a nie na wzorze, przy pięciu wartościach $w$ \dash{} bo ta
odpowiedniość zależy od konwencji, od tego, czy kara niesie połówkę, a
Program~\ref{prog:F04} zapisał, co się dzieje, gdy konwencję cytuje się z
pamięci.

Program~\ref{prog:P20} pracował przy $\lambda = \val{p26.p20.lambda}$. Jakie
to przekonanie?
\dotline
\nextframe
\end{fr}

\begin{fr}
\ans{Każda waga gaussowska wokół zera o szerokości $\val{p26.prior.tau}$}

\begin{center}
\begin{tabular}{rr}
\toprule
$\lambda$ & wynikająca szerokość a priori $\tau$ \\
\midrule
$\num{0.001}$ & $\val{p26.tau.0001}$ \\
$\num{0.01}$ & $\val{p26.tau.001}$ \\
$\val{p26.p20.lambda}$ & $\val{p26.tau.01}$ \\
$\num{1.0}$ & $\val{p26.tau.10}$ \\
\bottomrule
\end{tabular}
\end{center}

\textbf{Weight decay jest szerokością, a odczytanie go jako szerokości czyni
go sprawdzalnym.} Jeśli wytrenowane wagi warstwy mają rozrzut zupełnie
niepodobny do $\tau$, które wynika z twojego $\lambda$, to przekonanie i dane
się nie zgadzają, a to jest coś, na co można spojrzeć, a nie pokrętło, które
się kręci.

Nic w tej tabeli nie jest nową arytmetyką. To jeden odwrotny pierwiastek i
nigdy nie był drukowany obok $\lambda$.
\end{fr}

\begin{fr}
\begin{trapbox}
\textbf{\enquote{Weight decay to sztuczka, która akurat działa.}}

Jest dokładną konsekwencją jednego przekonania, a przekonanie da się
wypowiedzieć: wagi są gaussowskie wokół zera. Każdy inny rozkład a priori
daje inną karę \dash{} rozkład Laplace'a daje wartość bezwzględną zamiast
kwadratu, i dlatego $L_1$ produkuje zera, a $L_2$ nie, co jest faktem o
kształcie rozkładu w zerze, a nie faktem o optymalizatorach.

Powodem, dla którego czyta się to jak sztuczkę, jest to, że rozkładu a priori
nikt nie zapisuje. Jest tam, czy ktokolwiek go zapisze, czy nie, a jego
szerokość to $1/\sqrt{\lambda}$.
\end{trapbox}
\end{fr}

\begin{fr}
\begin{warning}
\textbf{MAP jest punktem i jest to prawdziwe ograniczenie.}

Wszystko powyżej maksymalizuje rozkład a posteriori, a potem go wyrzuca,
zachowując tylko argument maksimum. Ocena MAP nie umie więc powiedzieć, jak
bardzo jest pewna, nie umie powiedzieć, czy maksimum było ostre czy niemal
płaskie, i nie da się jej przenieść do dalszego rachunku.

Zachowanie całego rozkładu a posteriori to inny temat i należy do
Programu~\ref{prog:P28}: rozkład na parametrze zamiast punktu, z przedziałami
wiarogodnymi i arytmetyką, która czyni je obliczalnymi. Ten program kończy
tam, gdzie znaleziono maksimum.
\end{warning}

\mermaidfig{p26-prior-is-penalty}{Skąd bierze się kara, w jednym
kroku.}{wiara i logarytm}
\end{fr}

% ============================================================
\section{Score i dług z Części~VI}
\label{sec:P26-score}
\index{score}
\index{estymator!oparty na score}

\begin{fr}
Jedno wyprowadzenie zostało zaległe i nie jest ono własne temu programowi.

Program~\ref{prog:P21} spotkał wielkość będącą średnią po losowaniu i
potrzebował jej gradientu. Nazwał dwa sposoby jego uzyskania, zmierzył, co je
różni \dash{} trzy rzędy wielkości wariancji \dash{} a potem powiedział w
tylu właśnie słowach, że algebra za nimi należy do Programu~\ref{prog:P24}
albo do tego programu. Program~\ref{prog:P24} jej nie podał.

To cztery wiersze i jest to ten sam obiekt, który \S\ref{sec:P26-mle}
przyrównała do zera. Spójrz wstecz na tamtą pochodną: co było tą rzeczą,
której zero wyznaczało maksimum?
\dotline
\nextframe
\end{fr}

\begin{fr}
\ans{Pochodna logarytmu prawdopodobieństwa}

\[ s(x, \theta) \;=\; \frac{\partial}{\partial \theta}
   \ln p(x \mid \theta) \]

Nazywa się to \textbf{score}, a metoda największej wiarogodności jest
poleceniem \emph{przyrównaj sumaryczny score swoich danych do zera}. Tylko to
robiła \S\ref{sec:P26-mle}, pisząc $k/p - (n-k)/(1-p) = 0$.

Teraz pytanie o sam score. Uśrednij go po danych, o których model mówi, że je
zobaczysz \dash{} czyli weź jego wartość oczekiwaną pod
$p(x \mid \theta)$, przy tym samym $\theta$. Co wychodzi?
\dotline
\nextframe
\end{fr}

\begin{fr}
\ans{Zero, dokładnie}

Sprawdzone nad ułamkami przy $\val{p26.score.ps}$ różnych częstościach, bez
żadnej tolerancji.

Powód to jeden wiersz: $\Ex[s] = \sum_x p \cdot \frac{1}{p}
\frac{\partial p}{\partial \theta} = \frac{\partial}{\partial \theta}
\sum_x p = \frac{\partial}{\partial \theta} 1 = 0$. Prawdopodobieństwa sumują
się do jedynki, cokolwiek robi $\theta$, więc ich pochodna sumuje się do
niczego.

\textbf{Dlatego właśnie przyrównanie score do zera jest właściwym
poleceniem}, a nie tylko wygodą: przy prawdziwym parametrze score już średnio
wynosi zero, więc estymator, który zeruje score zaobserwowany, celuje we
właściwy warunek.
\end{fr}

\begin{fr}
Teraz dług. Chcesz gradientu $\Ex_{x \sim p_\theta}[f(x)]$, gdzie $\theta$
przesuwa rozkład, a $f$ jest czymkolwiek \dash{} nagrodą, symulatorem, oceną
człowieka.

Przeszkodą, którą nazwał Program~\ref{prog:P21}, jest to, że nie da się
zróżniczkować losowania. Ale wartość oczekiwana jest sumą po $x$ z
$p(x \mid \theta) f(x)$, a $p$ jest różniczkowalne, choć losowanie nie.

Zróżniczkuj tę sumę i zobacz, co kupuje ci score.
\dotline
\nextframe
\end{fr}

\begin{fr}
\ans{$\nabla_\theta \Ex[f] = \Ex\bigl[f(x)\,s(x,\theta)\bigr]$}

Bo $\nabla p = p \nabla \ln p$, więc
\[ \nabla_\theta \sum_x p\,f
   \;=\; \sum_x f\,\nabla_\theta p
   \;=\; \sum_x p\, f\, \nabla_\theta \ln p \]
a ostatnie wyrażenie jest wartością oczekiwaną, co znaczy, że da się je
oszacować przez losowanie.

\textbf{To jest estymator oparty na score z Programu~\ref{prog:P21} i są to
cztery wiersze metody największej wiarogodności.} Skrypt sprawdza to wobec
pochodnej wziętej z drugiej strony, dokładnie nad ułamkami, przy każdej
częstości i $\val{p26.score.cases}$ różnych funkcjach.

Jest to też dokładnie policy gradient. Pomnóż nagrodę, którą dostałeś, przez
gradient logarytmu prawdopodobieństwa akcji, którą podjąłeś, i uśrednij: to
wyrażenie nie jest przybliżeniem czegokolwiek, jest tożsamością, a jedyne, co
z nim nie tak, to wariancja, którą zmierzył Program~\ref{prog:P21}.
\end{fr}

\begin{fr}
Program zamyka się więc na jednym obiekcie widzianym cztery razy.

Score jest pochodną logarytmu prawdopodobieństwa. Przyrównanie go do zera to
metoda największej wiarogodności. Uśrednienie go daje zero, i dlatego jest to
właściwy warunek. Pomnożenie go przez nagrodę i uśrednienie różniczkuje przez
losowanie. A dodanie logarytmu rozkładu a priori do logarytmu wiarogodności
przed tym wszystkim to każda regularyzowana funkcja celu, jaka istnieje.

\mermaidfig{p26-set-it-to-zero}{Jedna pochodna i cztery rzeczy, które o nią
proszą.}{score, cztery uzycia}

Czego tu nie ma: co robić, gdy maksimum nie jest tym, na czym chcesz
poprzestać. Program~\ref{prog:P28} zachowuje cały rozkład a posteriori
zamiast jego szczytu, a Program~\ref{prog:P27} pyta, czy różnica między
dwoma ocenami jest większa niż szum w nich. Oba pytania ten program dopiero
uczynił możliwymi do zadania i żadne nie jest jego własne.
\end{fr}

% ============================================================
\begin{summarybox}
\sumitem{1--4}{\result{Estymator ma \textbf{dwa} sposoby bycia w błędzie.
  \emph{Obciążenie} to odległość średniej oceny od prawdy;
  \emph{wariancja} to odległość ocen od ich własnej średniej.} Program~\ref{prog:P21} dowiódł pierwszego o gradiencie z paczki i
  powiedział, że to nie wystarcza.}
\sumitem{5--6}{\result{\textbf{$\mathrm{MSE} = \text{obciążenie}^{2} + \Var$},
  dokładnie, sprawdzone przy dwudziestu jeden przyciąganiach nad ułamkami. Na
  omawianej populacji celowo obciążony estymator bije nieobciążony o
  $\val{p26.shrink.gain.pct}$ procent błędu średniokwadratowego \dash{}
  $\val{p26.mse.shrunk}$ wobec $\val{p26.mse.unbiased}$}.}
\sumitem{7--9}{\result{Zatem \enquote{nieobciążony} to nie \enquote{najlepszy}, a
  każda regularyzacja jest obciążeniem wziętym celowo. Optymalne przyciąganie
  zależy od nieznanej prawdy, i dlatego jest to spostrzeżenie, a nie
  przepis}.}
\sumitem{10--12}{\result{$n-1$ poprawia \textbf{dokładnie jedno obciążenie}:
  wariancja z próby z mianownikiem $n$ daje średnio $\frac{n-1}{n}$ wariancji
  populacji, wyliczone przy $n = 2$ do $5$, a mianownik $(n-1)$ daje ją
  dokładnie}.}
\sumitem{13--16}{\result{\textbf{Nie poprawia odchylenia standardowego.} Pierwiastek
  jest wklęsły, więc nierówność Jensena z Programu~\ref{prog:P19} sprawia, że
  poprawione $s$ zaniża \dash{} o $\val{p26.sd.short.2}$ procent przy
  $n = 2$ i wciąż o $\val{p26.sd.short.5}$ procent przy $n = 5$, czyli przy
  tym rozmiarze, od którego ludzie podają rozrzut}.}
\sumitem{18--20}{\result{\textbf{Metoda największej wiarogodności} wybiera parametr,
  który czyni zaobserwowane dane najbardziej prawdopodobnymi. Weź logarytm,
  żeby zamienić iloczyn w sumę (Program~\ref{prog:F03}) i uniknąć niedomiaru
  (Program~\ref{prog:P02}); maksimum się nie przesuwa, bo
  Program~\ref{prog:F05} mówi, że funkcja ściśle rosnąca zachowuje porządek}.}
\sumitem{21}{\result{Ocena bernoullowska to $k/n = \val{p26.bern.mle}$, wyprowadzona,
  a nie założona \dash{} \textbf{każda częstość policzona przez podzielenie
  sukcesów przez próby była oceną największej wiarogodności}}.}
\sumitem{22--23}{\result{A wariancja gaussowska wychodzi z $n$ w mianowniku,
  $\val{p26.gauss.mle.var}$ wobec $\val{p26.gauss.unbiased.var}$:
  \textbf{metoda największej wiarogodności jest obciążona}, dokładnie o
  czynnik z poprzedniej sekcji}.}
\sumitem{25--27}{\result{\textbf{Strata krzyżowoentropijna jest ujemnym logarytmem
  wiarogodności zaobserwowanych tokenów, na token.} Iloczyn
  prawdopodobieństw, logarytm, zmiana znaku, podzielenie przez ich liczbę
  \dash{} czyli definicja Programu~\ref{prog:P18} znak po znaku. Trenowanie
  modelu językowego to metoda największej wiarogodności}.}
\sumitem{28}{Trzy programy, jedno obliczenie:
  Program~\ref{prog:F02} podaje $\val{p26.f02.loss}$ nata, a
  Program~\ref{prog:F03} sekwencję $\val{p26.f03.tokens}$ tokenów;
  odtwarzają się nawzajem w obie strony, z powrotem do
  $\val{p26.back.loss}$ nata.}
\sumitem{30--31}{\result{\textbf{Label smoothing to metoda największej wiarogodności
  dla innego celu.} Przy $\varepsilon = \val{p26.ls.eps}$ na
  $\val{p26.ls.classes}$ klasach optimum to $\val{p26.ls.true}$ wobec
  $\val{p26.ls.other}$, czyli różnica logitów $\val{p26.ls.gap}$ \dash{}
  skończona, gdzie optimum celu one-hot jest nieosiągalne}.}
\sumitem{32--33}{\result{Ważenie klas szacuje parametry \emph{innej populacji}, a
  dzielenie przez liczbę tokenów jest tym, co czyni sekwencje różnej długości
  porównywalnymi}.}
\sumitem{34--36}{\result{\textbf{Kara jest logarytmem rozkładu a priori.} Pomnóż
  wiarogodność przez rozkład a priori, weź logarytm, zmień znak: logarytm a
  priori staje się składnikiem dodanym do straty, a gaussowski rozkład a
  priori o szerokości $\tau$ wnosi
  $\lVert\vect{w}\rVert^{2}/2\tau^{2}$}.}
\sumitem{37--38}{\result{Zatem $\lambda = 1/\tau^{2}$, dopasowane na gradientach,
  które dodaje Program~\ref{prog:P20}, a nie na wzorze.
  \textbf{Jego $\lambda = \val{p26.p20.lambda}$ to twierdzenie, że każda waga
  jest gaussowska o szerokości $\val{p26.prior.tau}$}, co da się sprawdzić na
  wytrenowanych wagach}.}
\sumitem{40}{\result{MAP jest punktem i nie umie powiedzieć, jak bardzo jest pewny.} Rozkład a posteriori jako rozkład należy do Programu~\ref{prog:P28}.}
\sumitem{41--43}{\result{\textbf{Score} to $\partial \ln p / \partial\theta$, a jego
  wartość oczekiwana wynosi dokładnie zero \dash{} bo prawdopodobieństwa
  sumują się do jedynki, cokolwiek robi $\theta$. Dlatego przyrównanie
  zaobserwowanego score do zera jest właściwym warunkiem}.}
\sumitem{44--45}{\result{A $\nabla \Ex[f] = \Ex[f\,s]$, dokładnie nad ułamkami przy
  każdej częstości i $\val{p26.score.cases}$ funkcjach: \textbf{estymator
  oparty na score, odłożony przez Program~\ref{prog:P21}, to cztery wiersze
  metody największej wiarogodności}, i jest tym, czym jest policy gradient}.}
\end{summarybox}

\canyou

\begin{testexercises}
\item Estymator ma obciążenie $\num{0.4}$ i wariancję $\num{0.09}$. Drugi ma
  obciążenie $0$ i wariancję $\num{0.30}$. Który jest średnio bliżej prawdy i
  o ile?
  \answerto{Błąd średniokwadratowy to $\text{obciążenie}^{2} + \Var$, więc
  pierwszy daje $\num{0.16} + \num{0.09} = \num{0.25}$, a drugi $\num{0.30}$.
  Obciążony jest bliżej, o $\num{0.05}$ \dash{} czyli teza sekcji w dwóch
  wierszach arytmetyki.}
\item Mierzysz wielkość pięć razy i podajesz średnią oraz odchylenie
  standardowe z próby z \code{ddof=1}. Podaj jedną rzecz, którą ta liczba nie
  jest, i powiedz, w którą stronę się myli.
  \answerto{Nie jest nieobciążoną oceną odchylenia standardowego. Poprawka
  czyni nieobciążoną \emph{wariancję}, a pierwiastek jest wklęsły, więc z
  nierówności Jensena podana wartość zaniża rozrzut \dash{} o jakieś
  $\val{p26.sd.short.5}$ procent przy $n = 5$ na populacji z tego programu, i
  zawsze w tę stronę, która sprawia, że wyniki wyglądają na powtarzalniejsze,
  niż są.}
\item Moneta daje $3$ orły w $10$ rzutach. Zapisz logarytm wiarogodności,
  znajdź ocenę największej wiarogodności i powiedz, dlaczego odpowiedź nie
  wymaga żadnego rachunku, gdy ma się poprzednie ramki pod ręką.
  \answerto{$3\ln p + 7\ln(1-p)$ z dokładnością do stałej, a jego pochodna
  $3/p - 7/(1-p)$ zeruje się przy $p = \num{0.3}$. Nie wymaga rachunku,
  bo bernoullowska ocena największej wiarogodności to $k/n$ dla każdego $k$ i
  $n$, co program wyprowadza raz.}
\item Przebieg treningowy podaje entropię krzyżową $\num{2.0}$ nata na token
  po $500$ tokenach. Podaj logarytm prawdopodobieństwa, jakie model
  przypisuje tej sekwencji, i powiedz, czy warto liczyć samo
  prawdopodobieństwo.
  \answerto{$-\num{2.0} \times 500 = -1000$ natów. Prawdopodobieństwo to
  $e^{-1000}$, czego nie utrzyma żaden format zmiennoprzecinkowy \dash{}
  czyli powód, dla którego Program~\ref{prog:P02} pracuje w logarytmach, i
  własny przykład Programu~\ref{prog:F03}, i dlatego raportuje się stratę, a
  nie wiarogodność.}
\item Dwa zespoły używają weight decay o sile $\num{0.01}$ i $\num{1.0}$ na
  tej samej architekturze. Powiedz, co każdy z nich twierdzi o wagach, i
  które twierdzenie jest mocniejsze.
  \answerto{Szerokości $\val{p26.tau.001}$ i $\val{p26.tau.10}$ wokół zera.
  Drugie jest o wiele mocniejsze \dash{} dziesięciokrotnie większe $\lambda$
  to rozkład a priori węższy o $\sqrt{10}$, więc twierdzi znacznie pewniej,
  że wagi są małe, i będzie znacznie wyraźniej w błędzie, jeśli nie są.}
\end{testexercises}

\begin{furtherproblems}
\item Pokaż, że dla estymatora przyciągniętego $c\bar{x}$ błąd
  średniokwadratowy jest minimalizowany przy
  $c = \mu^{2}/(\mu^{2} + \sigma^{2}/n)$, i powiedz, co się z nim dzieje, gdy
  $n$ rośnie.
  \answerto{Błąd to $\mu^{2}(c-1)^{2} + c^{2}\sigma^{2}/n$; zróżniczkuj i
  przyrównaj do zera. Gdy $n$ rośnie, $\sigma^{2}/n$ dąży do zera, a optimum
  dąży do $1$, więc przyciąganie jest warte tym mniej, im więcej masz danych
  \dash{} co jest tym samym stwierdzeniem co \enquote{mały model
  regularyzuj mocniej}.}
\item Score ma wartość oczekiwaną zero przy prawdziwym parametrze. Powiedz,
  czym to czyni sumaryczny score $n$ obserwacji, i użyj
  Programu~\ref{prog:P25}, by powiedzieć, jak się zachowuje, gdy $n$ rośnie.
  \answerto{Sumą $n$ niezależnych wielkości o średniej zero, więc ma średnią
  zero i wariancję $n$ razy większą od wariancji jednego składnika \dash{}
  pierwsza tożsamość Programu~\ref{prog:P25}. Podzielony przez $n$ koncentruje
  się na zerze w tempie $1/\sqrt{n}$, i dlatego ocena zerująca zaobserwowany
  średni score zbiega do prawdy.}
\item Rozkład a priori Laplace'a na każdej wadze jest proporcjonalny do
  $e^{-\lvert w\rvert/b}$. Podaj karę, której odpowiada, i powiedz w jednym
  zdaniu, co odróżnia go od przypadku gaussowskiego.
  \answerto{$\lVert\vect{w}\rVert_{1}/b$, czyli suma wartości bezwzględnych.
  Jego pochodna nie znika, gdy $w$ zbliża się do zera, a gaussowska znika,
  więc popycha do samego końca i produkuje dokładne zera \dash{} fakt o
  kształcie rozkładu a priori w zerze, a nie o optymalizatorze.}
\item Label smoothing przy $\varepsilon$ na $K$ klasach daje optymalną
  różnicę logitów
  \[ \ln\Bigl(\frac{1 - \varepsilon(K-1)/K}{\varepsilon/K}\Bigr) \]
  Powiedz, co się z nią dzieje, gdy $\varepsilon$ dąży do zera, i dlaczego
  takie zachowanie jest właściwe.
  \answerto{Dąży do nieskończoności, bo mianownik dąży do zera: bez
  wygładzania cel jest one-hot, a optimum nieosiągalne. Parametr wygładzania
  jest dokładnie tą rzeczą, która czyni optimum skończonym, więc granica
  odtwarzająca przypadek bez wygładzania jest sprawdzianem, że wzór jest
  dobry.}
\item Kolega zważył klasę rzadką w klasyfikatorze przez $10$ i raportuje, że
  trafność na tej klasie się poprawiła. Powiedz, co zostało oszacowane, i
  podaj jeden pomiar, który powiedziałby, czy to była właściwa rzecz do
  szacowania.
  \answerto{Parametry populacji, w której klasa rzadka jest dziesięć razy
  częstsza niż w danych. Czy o to chodzi, zależy od miksu przy wdrożeniu, więc
  pomiarem są częstości klas tam, gdzie model naprawdę pobiegnie \dash{}
  jeśli pasują do miksu przeważonego, wybór był uzasadniony, a jeśli pasują do
  oryginalnego, to wymieniono ogólną skuteczność na liczbę na slajdzie.}
\end{furtherproblems}
